\section*{Inner products}
\subsection*{From coordinates to a scalar product}
All vector spaces in the main text are real. An \textbf{inner product}
on \(V\) assigns a number \(\ip{u}{v}\) to each ordered pair, satisfying
\[
\ip{u}{v}=\ip{v}{u},\qquad
\ip{au+bw}{v}=a\ip{u}{v}+b\ip{w}{v},\qquad
\ip{v}{v}>0\quad(v\ne0).
\]
Linearity gives \(\ip{0}{v}=0\), and symmetry gives linearity in the
second argument. Define \(\norm{v}=\sqrt{\ip{v}{v}}\).
Vectors are orthogonal if their inner product is zero.
An orthonormal list consists of mutually orthogonal unit vectors.

The coordinate product on \(\R^n\) is one example. Two further examples
will keep the abstraction concrete.
On real \(m\times n\) matrices,
\[
\ip{A}{B}_F=\sum_{i,j}a_{ij}b_{ij}=\tr(A^\top B)
\]
is an inner product, with \(\norm{A}_F^2=\sum_{i,j}a_{ij}^2\).
Indeed the first two axioms follow from finite sums and the third from
the sum of squares. This is the \textbf{Frobenius inner product}.
On the real continuous functions on \([0,1]\),
\[
\ip{f}{g}=\int_0^1 f(t)g(t)\,dt
\]
is an inner product. Symmetry and linearity follow from integration.
If \(f\ne0\), continuity gives an interval of positive length on which
\(|f|\) is bounded below by a positive number, so \(\int_0^1 f^2>0\).
This argument uses continuity; a value at a single point would not suffice.

Expansion of the axioms gives
\[
\norm{u+v}^2=\norm{u}^2+2\ip{u}{v}+\norm{v}^2,\qquad
\norm{cv}=|c|\norm{v}.
\]
Thus Pythagoras holds for orthogonal vectors in any real inner-product
space, not only for coordinate vectors.

\begin{proposition}[Coefficients in an orthogonal list]
Nonzero mutually orthogonal vectors \(u_1,\ldots,u_r\) are independent.
If \(v=\sum_i c_i u_i\), then
\[
c_j=\frac{\ip{v}{u_j}}{\norm{u_j}^2}.
\]
For arbitrary \(v\), the vector
\[
p=\sum_{i=1}^r\frac{\ip{v}{u_i}}{\norm{u_i}^2}u_i
\]
satisfies \(v-p\perp u_j\) for every \(j\), hence \(v-p\) is perpendicular
to their whole span.
\end{proposition}
\begin{proof}
Taking the inner product of \(\sum_i c_i u_i\) with \(u_j\) leaves
\(c_j\norm{u_j}^2\). Applied to a zero combination, this proves that every
coefficient is zero. Applied to \(v\), it gives the formula.
For arbitrary \(v\), subtracting the displayed \(p\) gives
\[
\ip{v-p}{u_j}
=\ip{v}{u_j}-\frac{\ip{v}{u_j}}{\norm{u_j}^2}\norm{u_j}^2=0.
\]
Linearity gives orthogonality to any combination of the \(u_j\).
\end{proof}

\begin{theorem}[Projection, best approximation, and Bessel]
Let \(W=\operatorname{span}\{u_1,\ldots,u_r\}\) for a nonzero orthogonal
list. The vector \(p\) just defined is the unique vector of \(W\) with
\(v-p\perp W\). It is also the unique closest vector of \(W\) to \(v\).
For every \(w\in W\),
\[
\norm{v-w}^2=\norm{v-p}^2+\norm{p-w}^2.
\]
In particular,
\[
\sum_{i=1}^r\frac{\ip{v}{u_i}^2}{\norm{u_i}^2}\leqslant\norm{v}^2,
\]
with equality exactly when \(v\in W\).
\end{theorem}
\begin{proof}
The preceding proposition proves existence of the orthogonal residual.
If \(p'\in W\) has the same property, \(p-p'\) belongs to \(W\) and is
perpendicular to \(W\); its inner product with itself is zero, so \(p=p'\).
Now \(v-w=(v-p)+(p-w)\) is an orthogonal sum. Pythagoras proves the
distance identity. Its second term vanishes exactly at \(w=p\),
proving the minimizing assertion with uniqueness.
At \(w=0\), the identity gives
\(\norm{v}^2=\norm{v-p}^2+\norm{p}^2\).
Expansion in the orthogonal list yields
\(\norm{p}^2=\sum_i\ip{v}{u_i}^2/\norm{u_i}^2\).
Equality holds precisely when \(v=p\).
\end{proof}
For \(W=\{0\}\), use the empty list and \(p=0\).
For an orthonormal list \(q_i\), the coefficients simplify to
\(\ip{v}{q_i}\). When that list is a basis of \(V\), the Bessel inequality
is the equality \(\norm{v}^2=\sum_i\ip{v}{q_i}^2\).

\begin{proposition}[Schwarz and the triangle inequality]
For all \(u,v\),
\[
|\ip{u}{v}|\leqslant\norm{u}\norm{v},\qquad
\norm{u+v}\leqslant\norm{u}+\norm{v}.
\]
For \(v\ne0\), Schwarz equality means \(u\) is a multiple of \(v\);
for \(v=0\), equality always holds.
\end{proposition}
\begin{proof}
For \(v\ne0\), apply the preceding theorem to the one-vector list \(v\),
multiply by \(\norm{v}^2\), and take square roots. Its equality condition
gives the claim. If \(v=0\), both sides of Schwarz are zero.
Then
\[
\norm{u+v}^2
=\norm{u}^2+2\ip{u}{v}+\norm{v}^2
\leqslant(\norm{u}+\norm{v})^2.
\]
Both quantities whose squares are compared are nonnegative.
\end{proof}

\begin{example}[A weighted geometry]\label{eg:06:weighted}
Let
\[
G=\begin{pmatrix}3&-1\\-1&2\end{pmatrix},\qquad
\ip{x}{y}_G=x^\top Gy.
\]
Symmetry of \(G\) gives symmetry of this pairing, and products give
bilinearity. Its positivity follows by completing the square:
\[
x^\top Gx=3x_1^2-2x_1x_2+2x_2^2
=3(x_1-x_2/3)^2+\frac53x_2^2>0\quad(x\ne0).
\]
Such a symmetric matrix is called \textbf{positive definite}.
In this geometry \(\ip{e_1}{e_2}_G=-1\), so the coordinate axes are not
orthogonal. Projection of \(e_1\) onto \(\operatorname{span}\{e_2\}\) gives
\[
p=\frac{-1}{2}e_2,\qquad r=e_1-p=(1,1/2)^\top.
\]
Check \(\ip{r}{e_2}_G=-1+2(1/2)=0\). The squared norms are
\(\norm{e_1}_G^2=3,\ \norm{p}_G^2=1/2,\ \norm{r}_G^2=5/2\).
For a general candidate \(te_2\),
\[
\norm{e_1-te_2}_G^2=3+2t+2t^2
=\frac52+2(t+1/2)^2.
\]
The minimum and its unique location agree with the projection formula.
\end{example}

\section*{Orthogonal bases and complements}
\begin{theorem}[Gram--Schmidt]
For an independent list \(v_1,\ldots,v_k\), define
\[
u_1=v_1,\qquad
u_j=v_j-\sum_{i<j}\frac{\ip{v_j}{u_i}}{\norm{u_i}^2}u_i.
\]
Every \(u_j\) is nonzero, the list is orthogonal, and
\(\operatorname{span}\{u_1,\ldots,u_j\}
=\operatorname{span}\{v_1,\ldots,v_j\}\) at every stage.
Thus \(q_j=u_j/\norm{u_j}\) is an orthonormal basis of the original span.
\end{theorem}
\begin{proof}
Independence gives \(u_1=v_1\ne0\).
Suppose the conclusions hold before stage \(j\).
The projection calculation makes \(u_j\) perpendicular to every earlier
\(u_i\). If \(u_j=0\), then \(v_j\) is a combination of earlier \(u_i\),
and hence of earlier \(v_i\), contrary to independence.
The defining formula puts \(u_j\) in the span of \(v_j\) and earlier
vectors; rearranging puts \(v_j\) in the span of \(u_j\) and earlier
vectors. These two inclusions prove equality of the new spans.
Induction proves all claims. Normalizing preserves orthogonality and span.
\end{proof}

\begin{proposition}[Extension and decomposition]
An orthonormal list in a finite-dimensional inner-product space extends
to an orthonormal basis. If \(W\) is a subspace of such a space \(V\),
define
\[
W^\perp=\{v:\ip{v}{w}=0\text{ for all }w\in W\}.
\]
Then \(W^\perp\) is a subspace, and
\[
V=W\oplus W^\perp,\qquad
\dim W+\dim W^\perp=\dim V,\qquad
(W^\perp)^\perp=W.
\]
Here the direct sum means every vector has a unique decomposition into
one component in each subspace.
\end{proposition}
\begin{proof}
An orthonormal list is independent. Extend it to a basis and apply
Gram--Schmidt in that order; its initial orthonormal vectors are unchanged.
The remaining vectors normalize to give the claimed extension.

The defining equations of \(W^\perp\) are preserved by linear combinations.
Choose an orthonormal basis \(q_1,\ldots,q_r\) of \(W\) and extend it
to \(q_1,\ldots,q_n\) of \(V\). In
\(v=\sum_i\ip{v}{q_i}q_i\), the first \(r\) terms belong to \(W\),
and the rest are perpendicular to \(W\).
Conversely, if \(v\perp W\), its first \(r\) coefficients are zero.
Thus \(q_{r+1},\ldots,q_n\) form a basis of \(W^\perp\).
This proves existence, the dimension formula, and the double-complement
identity. If a vector belongs to both subspaces, it is perpendicular
to itself and is zero. Subtracting two decompositions therefore proves
uniqueness. Empty lists include the cases \(W=0\) and \(W=V\).
\end{proof}

\begin{example}[A plane: basis, orthogonalization, and projection]\label{eg:06:plane}
Consider \(W=\{(x,y,z)^\top:3x-2y+z=0\}\).
Because \(z=-3x+2y\), this space has dimension two. The vectors
\(a=(1,2,1)^\top\), \(b=(3,5,1)^\top\) lie in \(W\) and are independent:
the first and third coordinates of \(\alpha a+\beta b=0\) give
\(\alpha+3\beta=0,\alpha+\beta=0\), hence \(\alpha=\beta=0\).
Now
\[
a\cdot a=6,\quad a\cdot b=14,\quad
b-\frac{14}{6}a=\frac13(2,1,-4)^\top.
\]
With \(c=(2,1,-4)^\top\), we have \(a\cdot c=0,\ c\cdot c=21\).
Hence \(q_1=a/\sqrt6,\ q_2=c/\sqrt{21}\) is an orthonormal basis of \(W\).
Its complement has unit basis \(q_3=(3,-2,1)^\top/\sqrt{14}\).

For \(v=e_3\), the coefficients in the plane basis are
\(v\cdot q_1=1/\sqrt6\), \(v\cdot q_2=-4/\sqrt{21}\).
Thus
\[
p=\frac16a-\frac4{21}c
=\frac1{14}(-3,2,13)^\top,\qquad
r=e_3-p=\frac1{14}(3,-2,1)^\top.
\]
Directly, \(3p_1-2p_2+p_3=0\).
Also \(\norm{r}^2=1/14,\ \norm{p}^2=13/14\); their sum is \(\norm{e_3}^2\).
The closest-point identity is
\(\norm{e_3-w}^2=1/14+\norm{p-w}^2\) for every \(w\in W\).
\end{example}

\begin{proposition}[Projection matrices]
If \(Q\) has an orthonormal basis of \(W\subseteq\R^n\) as columns, then
\[
P_W=QQ^\top,\quad P_W^\top=P_W,\quad P_W^2=P_W,\quad
\operatorname{im}P_W=W,\quad \ker P_W=W^\perp.
\]
Also \(I-P_W=P_{W^\perp}\).
Conversely, a real symmetric idempotent matrix is the orthogonal
projection onto its image.
\end{proposition}
\begin{proof}
The coefficient formula is \(P_Wx=\sum_jq_j(q_j^\top x)=QQ^\top x\).
Since \(Q^\top Q=I\), transposing and squaring give the identities.
The decomposition proves the image, kernel and complementary formula.
Conversely, for \(P^\top=P=P^2\), \(Px\) is in its image and
\[
\ip{x-Px}{Py}=x^\top(P-P^\top P)y=0.
\]
Thus the decomposition relative to the image has component \(Px\),
as claimed. For the zero space the projector is zero.
\end{proof}
For the preceding plane, the formulas from its plane basis and its normal
\(n=(3,-2,1)^\top\) agree:
\[
\frac{aa^\top}{6}+\frac{cc^\top}{21}
=\frac1{14}\begin{pmatrix}5&6&-3\\6&10&2\\-3&2&13\end{pmatrix}
=I-\frac{nn^\top}{14}.
\]
The equality can be checked entry by entry; for example the \((1,2)\)
entry is \(2/6+2/21=6/14\).

For a matrix \(A\), \(Ax=0\) means \(x\) is perpendicular to every row.
Consequently
\[
\ker A=\operatorname{row}(A)^\perp,\qquad
(\ker A)^\perp=\operatorname{row}(A).
\]
These identities and rank--nullity give another proof that row rank
equals column rank: both equal \(n-\dim\ker A\).

\section*{Bilinear maps and matrices}
A map \(g:U\times V\to W\) is \textbf{bilinear} if it is linear in each
argument with the other fixed. It need not be linear in the pair:
for example \(g(u,v)=uv\) on \(\R\times\R\) is bilinear, but
\(g(2u,2v)=4uv\), not generally \(2uv\).
A scalar-valued bilinear map is a \textbf{bilinear form}.
Matrix multiplication is bilinear in its two matrix arguments;
the scalar-valued coordinate example is
\[
g_A(x,y)=x^\top Ay=\sum_{i=1}^m\sum_{j=1}^n a_{ij}x_i y_j,
\qquad A\in\R^{m\times n},\ x\in\R^m,\ y\in\R^n.
\]

\begin{theorem}[Representation of a bilinear form]
Given ordered bases \(B=(b_1,\ldots,b_m)\) of \(U\) and
\(C=(c_1,\ldots,c_n)\) of \(V\), every bilinear form \(g\) has a unique
matrix \(A\in\R^{m\times n}\) such that
\[
g(u,v)=[u]_B^\top A[v]_C,\qquad a_{ij}=g(b_i,c_j).
\]
This correspondence is a linear bijection between matrices and bilinear
forms. In particular the space of these forms has dimension \(mn\).
\end{theorem}
\begin{proof}
Write \(u=\sum_i x_i b_i\), \(v=\sum_j y_j c_j\).
Linearity in the first argument and then in the second gives
\(g(u,v)=\sum_{i,j}x_i y_j g(b_i,c_j)\), proving existence.
If two matrices give all the same values, substitute \(u=b_i,v=c_j\)
to recover each entry, proving uniqueness.
Conversely the coordinate formula for any \(A\) is bilinear.
Finally \(g_{A+D}=g_A+g_D\), \(g_{aA}=ag_A\), by the product laws.
Thus the correspondence is linear and bijective; the matrix units give
\(mn\) independent coordinate directions.
\end{proof}

\begin{proposition}[Symmetry and a change of basis]
On one space with the same basis in both positions, \(g_A\) is symmetric
if and only if \(A=A^\top\). A symmetric bilinear form is an inner product
if and only if \(x^\top Ax>0\) for every nonzero coordinate column \(x\).
If \(S\) converts new first-argument coordinates to old ones, and \(T\)
does the same for the second argument, the new matrix is
\[
A_{\rm new}=S^\top A_{\rm old}T.
\]
In particular the Gram matrix of an inner product changes by
\(G_{\rm new}=S^\top G_{\rm old}S\).
\end{proposition}
\begin{proof}
The scalar identity \(y^\top Ax=x^\top A^\top y\) proves the symmetry
claim by uniqueness of the representing matrix. Bilinearity is already
known, so precisely positivity remains for an inner product.
If the new coordinate columns are \(x',y'\), then the old ones are
\(Sx',Ty'\); substitution gives
\(g=x'^\top S^\top ATy'\), proving the transformation rule.
For an inner product, \(G_{ij}=\ip{b_i}{b_j}\); it is symmetric, and
\[
x^\top Gx=\left\|\sum_i x_i b_i\right\|^2>0\quad(x\ne0),
\]
because the basis is independent.
\end{proof}

\begin{example}[Reading and changing a bilinear form]\label{eg:06:bilinear}
For \(A=\begin{pmatrix}1&2\\3&-1\end{pmatrix}\),
\[
g_A(x,y)=x_1y_1+2x_1y_2+3x_2y_1-x_2y_2.
\]
For example \(g_A(e_1,e_2)=2\), whereas \(g_A(e_2,e_1)=3\);
this form is not symmetric.
Take the new basis \(b_1=(1,1)^\top,b_2=(1,-1)^\top\), with
\(S=\begin{pmatrix}1&1\\1&-1\end{pmatrix}\).
Then
\[
AS=\begin{pmatrix}3&-1\\2&4\end{pmatrix},\qquad
S^\top AS=\begin{pmatrix}5&3\\1&-5\end{pmatrix}.
\]
The \((1,2)\) entry is checked directly:
\(g_A(b_1,b_2)=(1,1)(-1,4)^\top=3\).
For the ordinary dot product the same basis has Gram matrix
\(S^\top S=2I\). Thus the basis is orthogonal with squared lengths two,
and its coordinate norm formula is \(2x_1'^2+2x_2'^2\).
\end{example}

\paragraph{Calculus connection.}
For a twice continuously differentiable real function \(f\) near a point,
its Hessian \(H=(\partial^2f/\partial x_i\partial x_j)\) represents the
bilinear form \((u,v)\mapsto u^\top Hv\). Equality of mixed partials under
this regularity assumption makes it symmetric. Symmetry alone does not
imply positivity: \(f(x_1,x_2)=x_1^2-x_2^2\) has Hessian
\(\operatorname{diag}(2,-2)\).
